\section*{Part IV: How to Discover Willingness to Pay}

\subsection*{24. Start with Revealed Sacrifice, Not an Invented Maximum}

People can answer ``What would you pay?'' without experiencing the consequences of their reply. They may want to be supportive, seem sophisticated, signal frugality, or imagine a different offer. Begin instead with the last time they made a comparable sacrifice. Which invoices, purchases, travel, workarounds, training expenditures, hours, and postponed projects are real? What did they reject and why?

A developer who purchased a ₹6,000 specialist review last month demonstrates a different kind of evidence from someone who says expert mentoring is worth ₹1 lakh. Neither observation fixes the WTP for your proposed program: the scope, trust, urgency, and alternatives differ. But actual past expenditure anchors the discussion in behavior and reveals budget categories.

Ask for a chronology. Avoid requests for sensitive records unless necessary and freely consented to. A good interview might end with a willingness-to-pay hypothesis; it should not begin with a demand for a number.

\subsection*{25. Discovery Interviews Should Reconstruct Decisions}

The strongest interviews are investigations of behavior. Ask: What triggered the search? What was going wrong? What was the cost of leaving it alone? What did you try 1st? Who else was involved? What options reached the shortlist? Which were too risky, too time-consuming, or hard to justify? What did you actually pay? What happened afterward? What would you do again?

Then test the counterfactual: if our proposed solution did not exist, what would you do? Ask for concrete evidence of the alternative. The answer may change as you describe a prototype, so distinguish the original behavior from the later reaction to your pitch. Record quotations separately from your interpretations.

A useful interview sample includes recent purchasers, nonpurchasers, buyers who chose a competitor, and people who abandoned the problem. If you interview only people who agreed to a founder meeting, you may systematically oversample the helpful and curious. Research should challenge the founder's expectation, not produce a room full of encouragement.

\subsection*{26. The Evidence Ladder: What Each Signal Proves}

Not all customer feedback carries identical evidentiary weight. The evidence ladder arranges signals from weakest to most conclusive:

\begin{table*}[t]
\centering
\small
\caption{The Evidence Ladder: Signals of Commercial Demand and Their Epistemic Limits}
\begin{tabularx}{\textwidth}{p{0.24\textwidth}p{0.36\textwidth}X}
\toprule
\textbf{Observed Signal} & \textbf{Strongest Defensible Inference} & \textbf{Important Epistemic Limitation} \\
\midrule
Compliment or positive rating & Comprehension, courtesy, or intellectual interest & Proves nothing about willingness to sacrifice scarce money \\
Previous purchase for same job & Established category budget and willingness to pay & Prior offering or market conditions may differ \\
Concrete cost of workaround & Problem is recognized as consequential & Modeled labor costs may be inflated or non-redeemable \\
Forced trade-off among offers & Relative attribute prioritization & Subject to hypothetical bias and laboratory distortion \\
Priced landing-page response & Click-through interest under visible price terms & No money changes hands; low friction to express interest \\
Refundable deposit & Material commitment of liquid funds & Refund guarantee drastically reduces perceived downside \\
Paid pilot under clear terms & Real exchange in that specific setting & Small sample does not establish organizational repeatability \\
Full-price sale with low regret & Genuine exchange of value experienced & Early adopters may be unusually motivated or forgiving \\
Repeat acquisition at scale & Durable commercial pattern under normal conditions & Competitors, macroeconomics, and channels shift over time \\
\bottomrule
\end{tabularx}
\end{table*}

A paid transaction is stronger evidence than an interview, but 1 enthusiastic buyer is not a market. A 90 percent response from 6 handpicked friends is not a scalable conversion rate. Every observation is conditional on a selection process, and the founder must understand that process before extrapolating.

\subsection*{27. Van Westendorp: A Map of Price Perceptions}

The Van Westendorp technique asks respondents when an offer seems too cheap to trust, inexpensive, expensive but plausible, and too expensive. It can reveal how a defined population interprets prices within a familiar category. It does not directly estimate a true revenue-maximizing price, buyer surplus, or paid conversion. If respondents imagine different levels of service, their answers are not even about the same object.

Use it to expose reference frames and potential price ranges after specifying the product. Segment the data by job and buyer context. Do not call the intersection of its lines ``the market's mathematically optimal price.'' It is a perception survey, and survey answers may not survive actual purchasing friction.

\subsection*{28. Gabor-Granger and Discrete-Choice Experiments}

Gabor-Granger approaches test hypothetical purchase likelihood at explicit prices; discrete-choice and conjoint designs vary attributes including price so that respondents must choose among bundles. These can help the founder discover whether expert review matters more than extra video content, whether shorter duration increases preference, and how buyers trade cost for risk reduction.

Their usefulness depends on credible alternatives, representative respondents, sensible attribute combinations, and attention to hypothetical bias. Asking a respondent to choose among 3 imaginary learning programs that none can inspect may produce elaborate but weak data. For a tiny early venture, a few well-executed real offers can be more informative than a complex statistical model built on an unrepresentative convenience sample.

Structured choice research becomes stronger when validated against paid behavior. It can guide what to test, not confer certainty about what will sell.
